\chapter{二次型；初等几何.}

以现代面目出现的二次型理论，其历史很难追溯到 18 世纪下半叶之前，而且我们将会看到，它的发展主要是为了回应数论、分析和力学的需要。但这一理论的基本观念实际上从“欧几里得”几何的开端起就已经出现，并且构成了这种几何的框架。因此，若不至少以概述的方式谈谈自古以来“初等几何”的发展，便无法追溯它的历史。当然，我们只能勾勒少数几个一般观念的演变，读者不应指望在这里找到关于这个或那个具体定理的历史的精确资料，对此只需求助于专门的历史著作或教学著作即可。\(^{1}\)同样不言而喻的是，当我们在下文谈到同一定理用各种代数语言或几何语言表达时的各种可能解释时，我们完全无意宣称这些“翻译”在各个时代都像今天这样为人们所熟悉；恰恰相反，本篇笔记的主要目标正是要表明，数学家们是如何极其缓慢地逐渐意识到这些外表往往迥然不同的问题之间的联系；我们还必须表明，他们在这样做时，如何被引导着给古人遗留下来的一大堆几何定理注入某种连贯性，并最终试图精确地划定“几何”一词究竟应当理解为什么。

如果撇开巴比伦人关于二次方程求解公式的发现（{[}232{]}，第 183–189 页）不谈，那么就应当在几何外衣之下去寻找二次型理论诸主要概念的诞生。这些概念最初以距离的平方（在平面上或三维空间中）的形式出现，与之对应的“正交性”概念则借助直角引入，欧几里得把直角定义为平角的一半（《几何原本》第 I 卷，定义 10）；距离与直角这两个概念由毕达哥拉斯定理——欧几里得大厦的拱顶石——联系在一起。\(^{2}\)角的概念在希腊数学中似乎很早就被引入了（这无疑是得自巴比伦人，后者凭借长期的天文经验已经习惯于使用角）。众所周知，在古典时期，只有小于两直角的角才有定义（欧几里得的这个“定义”与他给出的直线和平面的定义一样，既含混又不可用）；方向的概念并未被提炼出来，尽管欧几里得使用了（既无公理亦无定义）直线把平面分成两个区域这一事实，而且在必要之处他还会仔细地区分这两个区域。\(^{3}\)在这个阶段，平面旋转群的观念只能以非常不完善的方式显露出来，即通过半线的非定向角的加法（它同样是欧几里得未加解释地引入的），而这种加法原则上只在和至多等于两直角时才有定义。\(^{4}\)至于三角学，它受到几何学家的轻视，被交给土地测量者和天文学家；正是后者（阿里斯塔克斯、喜帕恰斯，尤其是托勒密{[}255{]}）建立了直角三角形（平面的或球面的）边与角之间的基本关系，并编制了最初的数表（这些表给出半径为 r 的圆上被角 \(\theta < \pi\) 截出的弧所对的弦，换句话说就是数 \(2r \sin \frac{\theta}{2}\)；更便于使用的正弦的引入应归功于中世纪的印度数学家）；在计算这些表时，当时尚不为人知的弧加法公式被换成对托勒密定理的等价使用（它也许可以追溯到喜帕恰斯）

关于圆内接四边形的定理。还应当指出，欧几里得和海伦就任意平面三角形的边与角给出了与公式

\[
a ^ {2} = b ^ {2} + c ^ {2} - 2 b c \cos A
\]

等价的命题；但由于缺乏向量演算的观念——它要到 19 世纪才会出现——人们很难在其中看出与度量型相伴的双线性型概念的初次现身。

位移（或者说运动，这两个概念的区别在古代——甚至在很久之后——都不清楚）已为欧几里得所知；但出于我们无从知晓的原因，他似乎对使用位移怀有明显的反感（例如在“三角形全等的情形”中，人们得到的印象是，他之所以使用位移概念，只是因为他未能表述出一条适当的公理（{[}153 e{]}，第 I 卷，第 225–227 页及第 249 页））；然而，为了定义回转锥面和球面（《几何原本》第 XI 卷，定义 14 和 18），他求助于的恰恰是位移（绕一条轴的旋转）概念，阿基米德定义回转二次曲面时用的也是这一概念。但是，作用于整个空间的变换这一一般观念，在 18 世纪末之前几乎与数学思想无缘；⁵而在 17 世纪之前，人们既找不到运动合成概念的任何痕迹，更谈不上位移的合成。当然，这并不是说希腊人没有特别意识到图形的“正则性”和“对称性”——我们现在把这些概念与位移群的概念联系在一起：他们关于正多边形、尤其是正多面体——其全部数学中最出色的篇章之一——的理论，正好证明情况相反。⁶

最后，就我们所关心的领域而言，希腊数学最后一项本质贡献是圆锥曲线的理论（至于二次曲面，希腊人只知道某些回转二次曲面，而且除球面外，他们对这些曲面的研究都没有走得很远）。这里值得注意的是，尽管希腊人从未产生解析几何基本原理的想法（主要因为缺少一种便于使用的代数），他们在研究特殊“图形”时通常使用相对于平面上两条（甚至多于两条）轴的“纵标”（这些轴与图形紧密相连，这正是他们的方法与费马和笛卡尔的方法的根本差异之一，后者的轴是独立于所考虑的图形而固定的）。特别地，在倍立方问题中出现的第一批圆锥曲线（圆除外）的例子，是由方程 \(y^{2}=ax, y=bx^{2}, xy=c\) 给出的曲线（欧多克索斯的学生梅内克缪斯，公元前 4 世纪中叶）；\(^{7}\)而在研究与这些曲线有关的问题时，最常使用的是圆锥曲线的方程（通常相对于由一条直径以及曲线与它的一个接触点处的切线所组成的两条斜轴）（而“焦点”性质只起着非常不显眼的作用，这与仅仅上溯到 19 世纪的渊博传统会使我们相信的恰好相反）。在这门浩繁的理论中，我们在这里尤其应当记住共轭直径的概念（阿基米德已经知道它），以及今天用来定义一点的极线的那条性质，它由阿波罗尼奥斯{[}153 b{]}就该点在圆锥曲线之外的情形给出（因此对他而言，极线就是由从该点引出的各切线的切点连成的直线）；从我们的观点看，这是关于异于度量型的二次型的“正交性”的两个例子，但可以想见，这些概念与经典的垂线概念之间的联系，在那个时代是不可能被设想到的。

在笛卡尔和费马之前，几乎没有别的进展可以指出；但随着解析几何的开端，二次型的代数理论开始从它的几何外壳中脱胎而出：费马知道平面上的二次方程表示一条圆锥曲线（{[}109{]}，第 I 卷，第 100–102 页；法文译本，第 III 卷，第 84–101 页），并对二次曲面勾画了类似的想法（{[}109{]}，第 I 卷，第 111–117 页；法文译本，第 III 卷，第 102–108 页）。随着 18 世纪中二维和三维解析几何的发展，该理论的两个中心问题出现了（尤其是联系着圆锥曲线和二次曲面）：把二次型化为平方和，以及寻求它关于度量型的“轴”。对圆锥曲线来说，这两个问题太初等，引不起重要的代数进展；对任意多个变量的情形，第一个问题由拉格朗日于 1759 年解决，与多元函数的极大值有关（{[}191{]}，第 I 卷，第 3–20 页）。但这个问题几乎立刻被求轴的问题所掩盖，甚至早在秩的不变性被表述之前就是如此；\(^{8}\)至于惯性定律，它直到 1850 年前后才由雅可比（{[}171{]}，第 III 卷，第 593–598 页）发现，他以与今天相同的论证证明了它，而西尔维斯特（{[}304{]}，第 I 卷，第 378–381 页）则满足于把它当作近乎显然的事实加以陈述。\(^{9}\)

把二次曲面化为它的轴的问题，已经呈现出比圆锥曲线的相应问题大得多的代数困难；最先着手处理这个问题的欧拉，尚无力证明特征值的实性，他在一次缺乏说服力的论证尝试之后便承认了这一点（{[}108 a{]}，(1)，第 IX 卷，第 379–392 页）。\(^{10}\)这一点大约在 1800 年{[}140{]}才被正确建立，而关于任意 n 个变量的型的相应定理的证明，则必须等到柯西（{[}56 a{]}，(2)，第 IX 卷，第 174–195 页）。也是柯西，在大约同一时间，证明了给出特征值的特征方程在一切正交换轴之下不变（{[}56 a{]}，(2)，第 V 卷，第 252 页）；\(^{11}\)但对 n = 2 或 n = 3，这种不变性凭特征值由相应的圆锥曲线或二次曲面的轴给出的几何解释，在直观上是“显然的”。此外，在关于这个课题的研究过程中，特征值的初等对称函数也自然地出现过（带有种种几何解释，特别是与阿波罗尼奥斯关于共轭直径的定理相联系），尤其是判别式，它（在 n = 2 的情形，联系着二次方程理论而早已为人所知）在 n = 3 的情形由欧拉首次引入（{[}108 a{]}，(1)，第 IX 卷，第 382 页）；后者是在给二次曲面分类时遇到它的（他借此表达二次曲面没有无穷远点的条件），却未提及它关于正交换轴的不变性。但稍后，随着整系数二次型的算术理论的开端，拉格朗日（就 n = 2）注意到了判别式在线性但非正交的变量代换之下不变性的一个特例（{[}191{]}，第 III 卷，第 699 页），高斯则就 n = 3 建立了判别式对一切线性变换的“协变性”（{[}124 a{]}，第 I 卷，第 301–302 页）。\(^{12}\)一旦柯西和比内证明了行列式乘法的一般公式，把高斯的公式推广到任意多个变量的情形便是直接的了；正是这一点，在大约 1845 年，给了一般的不变量理论以最初的推动。

在希腊人那里取代位移理论的两个概念——运动的概念和图形“对称”的概念——之外，17 和 18 世纪又加上了第三个，即直角坐标轴变换的问题，它与位移理论实质上等价。欧拉为这个问题撰写了好几部著作，他最看重的是为换轴公式求得便于使用的参数表示。力学日后对他为 n = .3 的情形为此目的引入的那三个角作了什么用处，这是人所共知的（{[}108 a{]}，(1)，第 IX 卷，第 371–378 页）。但他并不止步于此：他在 1770 年考虑了任意 n 的正交变换的一般问题，指出这里通过引入 \(n(n - 1)/2\) 个角作为参数即可达到目的，最后，对 n = 3 和 n = 4，他给出了旋转的有理表示（分别是 4 个齐次参数、以及由一个关系式联系起来的 8 个齐次参数的函数），它们正是后来借助四元数理论得到的那些表示，而他并未指出其来历（{[}108 a{]}，(1)，第 VI 卷，第 287–315 页）。\(^{13}\)

另一方面，欧拉还指出了如何解析地表达对平面图形“对称”的寻求，正是为了这个目的，他被引向实质性地证明：一个平面位移是一个旋转，或是一个平移，或是一个平移后面接一个对称（{[}108 a{]}，(1)，第 IX 卷，第 197–199 页）。力学的兴起此时还导致对位移的一般研究：但起初涉及的只是与连续运动相切的“无穷小”位移：在托里拆利、罗贝瓦尔和笛卡尔关于运动合成以及平面运动的瞬时旋转中心的研究中，出现的显然只是这种位移（参见第 177 页）。后者由约翰·伯努利给出了一般的定义；达朗贝尔于 1749 年、欧拉于次年推广了这个概念，证明了对使一点固定不动的运动存在瞬时旋转轴。有限位移的相应定理直到 1775 年才由欧拉{[}108 b{]}陈述出来，在他于此同时发现旋转的行列式等于 1 的那篇论文中；次年，他证明了平面相似变换存在不动点（{[}108 a{]}，(1)，第 XXVI 卷，第 276–285 页）。但要最终得到关于有限位移和无穷小位移的一个连贯理论，还须期待沙勒始于 1830 年的工作{[}60 a{]}。

这样我们就来到了可称之为几何学黄金时代的时期，它大致处于蒙日《画法几何》（1795 年）{[}225{]}的出版与 F. 克莱因的“埃尔朗根纲领”（1872 年）（{[}182{]}，第 I 卷，第 460–497 页）这两个日期之间。这一几何学的骤然复兴带给我们的本质进展如下：

\begin{enumerate}
\def\labelenumi{\Alph{enumi})}
\item
  无穷远元素（点、直线或平面）的概念由德萨格在 17 世纪引入{[}84{]}，但在 18 世纪里，它除了被滥用为一种说法之外几乎没有得到表达；彭赛列{[}252{]}恢复了它的地位并系统地使用它，从而把射影空间弄成了一切几何现象的一般框架。
\item
  与此同时，随着蒙日、尤其是彭赛列的工作，向复射影几何的过渡得以完成。在 18 世纪中只被零星使用的虚点概念，在这里（与无穷远点的概念同时）得到利用，以给出独立于实仿射几何的“图形情形”的命题。如果说最初为支持这些创新而提出的辩护还很笨拙（尤其是出自“综合”几何学派成员之手的辩护，在这个学派中，坐标的使用最终被视为一件不体面的事），那么不容错过的是，在那里可以辨认出，以蒙日的“偶然关系原理”或彭赛列的“连续性原理”的名义出现的，正是现代代数几何“特殊化”思想的最初萌芽。\(^{14}\)
\end{enumerate}

从这些概念得出的最早结果之一，是注意到在复射影空间中，所有非退化的圆锥曲线（相应地，二次曲面）都属于同一类；这引导彭赛列发现了“迷向”元素：“平面上任意放置的各个圆”，他说，“因此并不是彼此完全独立的，不像初看时所可能认为的那样，它们在理想意义下于无穷远处有两个公共的虚点”（{[}252{]}，第 I 卷，第 48 页）。后来，他又以同样的方式引入了“脐形曲线”，即所有球面公共的无穷远虚圆锥曲线（{[}252{]}，第 I 卷，第 370 页）；即使他没有特别谈及球面的迷向母线，他至少明确强调了所有二次曲面都存在实的或虚的直母线（同前引文，第 371 页）；\(^{15}\)这些概念，他的后继者们（尤其是

普吕克和沙勒）对它们的使用比他本人还要多，特别是在圆锥曲线和二次曲面的“焦点”性质的研究中。

\begin{enumerate}
\def\labelenumi{\Alph{enumi})}
\setcounter{enumi}{2}
\item
  点变换和变换合成的概念也同样得到了一般的表述，并被系统地引进作为证明手段。除了位移和投影之外，直到那时人们只知道一些特殊的变换：拉伊尔和牛顿使用过的 \(x' = a / x\)、\(y' = y / x\) 型的某些平面射影变换，克莱罗和欧拉的“仿射变换”\(x' = ax\)、\(y' = by\)（{[}108 a{]}，(1)，第 IX 卷，第 XVIII 章），最后还有再次由牛顿以及麦克劳林和布雷肯里奇给出的若干特殊二次变换。蒙日在其《画法几何》中展示了平面投影在三维几何中能得到多么充分的利用。在彭赛列那里，一种已被用到泛滥的系统化证明程序，就是通过投影把圆锥曲线的性质化归为圆的性质（笛卡尔和帕斯卡已经偶尔使用过这种方法）；而为了从二次曲面过渡到球面，他发明了空间的射影变换的第一个例子，即“透射”（{[}252{]}，第 I 卷，第 357 页）；最后，也是他引入了曲线到自身的双有理变换的最初例子。1827 年，默比乌斯（{[}223{]}，第 I 卷，第 217 页）（以及 1830 年沙勒的独立工作（{[}60 b{]}，第 695 页））定义了最一般的线性射影变换；与此同时出现了反演以及其他类型的二次变换，对它们的研究将开创双有理变换的理论，这一理论将在 19 世纪下半叶发展起来。
\item
  对偶的概念完全公开登场，并且被发现自觉地与双线性型理论联系在一起。关于圆锥曲线的极点与极线的理论，自阿波罗尼奥斯以来只在德萨格和拉伊尔那里取得过一些进展，蒙日把它推广到了二次曲面，他和他的学生们都看到了借此把已知定理变换成新结果的可能性。\(^{16}\)但是，把这些评述建构成“互反配极”变换理论中的一般方法、并把它变成一种特别有效的发现工具，这一功劳仍然要归于彭赛列。稍后，特别是在热尔岗、普吕克、默比乌斯和沙勒那里，对偶的一般概念摆脱了它与二次型的联系（这种联系在彭赛列那里还过于紧密）。特别地，默比乌斯在考察三维空间中对偶（由一个双线性型定义）的各种可能性时，于 1833 年发现了关于交错双线性型的对偶（{[}223{]}，第 I 卷，第 489–515 页），\(^{17}\)它在 19 世纪尤其以“线性线丛”理论的形式得到研究，并联系着凯莱、格拉斯曼和普吕克约在 1860 年引入的“线几何学”和“普吕克坐标”而发展起来。
\item
  从射影几何的开端起，对经典几何的各种性质及其与射影空间的关系的深入研究，很快导致把它们划分为“射影性质”和“度量性质”；把这一区分视为日后成为现代结构概念的那种东西在当时最清晰的表现之一，无疑并不算夸张。但是最先引入这一区分及其术语的彭赛列，已经意识到联系着这两类性质的东西；在《论著》中处理关于角的问题时，这些性质“似乎不属于我们称之为射影的那类性质 \ldots，然而它们毕竟以如此简单的方式”，他说，“得自构成{[}本书{]}基础的那些原理 \ldots，以致我不认为有任何别的几何理论能够以既更直接又更简单的方式导出它们。如果考虑到图形的射影性质必然是图形所能具有的性质中最一般的性质，这就不会令人惊讶；这样一来，它们必定像简单的推论一样，把该领域所有其他的性质或特殊关系都包括进来”（{[}252{]}，第 I 卷，第 248 页）。说实话，在这番宣言之后，看到他以非常迂回的方式处理关于角的问题——把角与圆锥曲线的焦点性质联系起来，而不是直接引入圆点——未免有点令人惊讶；而事实上，直到 30 年后，拉盖尔（当时还是综合理工学校的学生）才借助这些直线与同起点的各迷向直线的交比给出直角的表达式（{[}192{]}，第 II 卷，第 13 页）。最后，随着凯莱（{[}58{]}，第 II 卷，第 561–592 页）的工作，平面图形的“度量”性质无非是该图形添上圆点之后的“射影”性质这一基本思想得到了明确的表达——这是通向“埃尔朗根纲领”的一个决定性中途站。
\item
  将在 1830 年前后问世的双曲型非欧几何，起初多少置身于我们正勾勒其主要脉络的那场运动之外。这门新几何产生于本质上属于逻辑层次、涉及经典几何基础的关切，其发明者\(^{18}\)以与欧几里得几何相同的公理化“综合”形式表述它，而且它与射影几何毫无联系（在这门几何中唯一平行的概念消失了，因此按照经典模式引入射影几何甚至似乎先验地就被排除）；无疑正是由于这个缘故，在很长时间里，它几乎引不起法国、德国或英国的射影几何学派的任何兴趣。于是，当凯莱在前引的那篇基本论文（{[}58{]}，第 II 卷，第 561–
\end{enumerate}

\begin{enumerate}
\def\labelenumi{\arabic{enumi})}
\setcounter{enumi}{591}
\tightlist
\item
  592 页）中想到用一条任意圆锥曲线（他称之为“绝对形”）来代替圆点（被看作“相切退化”的圆锥曲线）时，他根本没有想到把这个想法与罗巴切夫斯基–鲍耶的几何联系起来，尽管他指出了他的构想如何导出两点间“距离”的新表达式，也提到了它与球面几何的联系。局势大约在 1870 年发生变化，此时非欧几何随着罗巴切夫斯基著作的传播、高斯著作的出版以及黎曼的就职演讲而来到数学新闻的前沿。沿着黎曼开拓的道路，贝尔特拉米在不知道凯莱工作的情况下，于 1868 年重新找到了后者给出的距离表达式，但处在相当不同的语境中：他把一个圆的内部看作一张常曲率曲面的像，其中测地线由直线来表示{[}18 a{]}；两年后，克莱因（独立于贝尔特拉米）对这些不同的观点作了综合，并通过发现椭圆型非欧空间（{[}182{]}，第 I 卷，第 254–305 页）而使之完备。\(^{19}\)
\end{enumerate}

\begin{enumerate}
\def\labelenumi{\Alph{enumi})}
\setcounter{enumi}{6}
\tightlist
\item
  在我们此处考察的时期的后半段，还插入了一个批判性反思的阶段，在此期间，“综合”几何的支持者们不满足于把坐标从证明中驱逐出去，甚至打算在几何公理中也弃用实数。这一学派的主要代表是冯·施陶特，他基本上成功地完成了这一杰作{[}325{]}，它在其时代、乃至进入 20 世纪很久之后都备受推崇；虽然今天人们对这种层次的思想不再给予同样的重视——其富有成效的应用前景被证明相当有限——但仍应承认，冯·施陶特及其弟子们的努力有助于澄清实的或复的“标量”在经典几何中的作用，并恰恰以此引入了在任意基域上的几何这一现代观念。
\end{enumerate}

大约在 1860 年，“综合”几何达到了它的顶峰，但其统治的终结也在迅速临近。在整个 18 世纪一直笨重而粗陋的解析几何，在拉梅、博比利耶、柯西、普吕克和默比乌斯手中，终于获得了使它得以与对手平起平坐的优雅和简洁。尤其是，大约 1850 年以后，群和不变量的观念终于获得了精确的表述，一点点地占领了舞台，人们看到，经典几何的定理无非是相似群的不变量或共变式之间的恒等关系的表达，\(^{20}\)正像射影几何的定理表达射影群的共变式之间的恒等关系（即“合冲”）一样。正是 F. 克莱因在著名的“埃尔朗根纲领”（{[}182{]}，第 I 卷，第 460–497 页）中权威地陈述了这一论点，他在其中主张抛弃“综合”倾向与“解析”倾向之间毫无结果的争论；他说，如果说针对后者让任意的坐标系扮演特权角色所提出的指责，“就先前使用坐标方法时的缺陷而言，在大多数情况下是正当的，那么当问题在于对这种方法的合理运用时，它就站不住脚了……空间直观的领域并没有被解析方法所禁止……”，他还强调，“人们不应低估一个合用的格式给后续研究带来的好处，因为它可以说走在思想的前面”（同前引文，第 488–490 页）。

于是人们最终得到了按照定理所从出的群对“几何”定理所作的一种合理的“结构性”分类：射影几何对应线性群，度量问题对应正交群，“线性线丛”几何对应辛群。但在这道无情的强光之下，经典几何——代数几何和微分几何除外，\(^{21}\)它们从此被确立为自治的学科——骤然熄灭，失去了全部光彩。基于变换之使用的诸方法的推广，早已使新定理的产出变得有些机械：沙勒 1837 年在其《历史概览》中说：“今天，人人都可以挺身而出，任取一条已知的真理，把它交付给各种一般的变换原理；他将得出别的真理，不同的或更一般的；而这些又可接受类似的操作；于是，从第一条真理推导出来的新真理的数目几乎可以倍增到无穷 \ldots{} 因此，在科学的现状下，谁愿意，谁就能够加以推广并在几何中创造；为这座大厦添上一块石头，天才已不再是不可或缺的”（{[}60 b{]}，第 268–269 页）。但随着不变量理论的进展，局势变得清楚多了：它终于（至少对“经典”群而言）表述出了一些一般方法，原则上能以纯粹自动的方式写出一切代数共变式及其全部“合冲”；这是一场胜利，它同时宣告了经典不变量理论本身以及“初等”几何作为研究领域的死亡，\(^{22}\)后者在实践中已变成它的一本简单词典。无疑，在可以随意大量炮制出来的无穷多条“定理”当中，没有什么办法先验地预见哪些定理用恰当的几何语言陈述出来会具有可与经典结果相比拟的简洁与优美，那里还留有一块狭窄的领域，众多业余爱好者仍乐于在其中耕耘（三角形、四面体的几何，低次代数曲线与曲面的几何，等等）。但对职业数学家来说，矿藏已经采竭，因为那里不再有任何可能波及数学其他部分的结构性问题；群论与不变量理论的这一篇章，在另行通知之前，可以视为已经终结。\(^{23}\)

因此，在埃尔朗根纲领之后，欧氏几何与非欧几何从纯代数的观点看，变成了或多少便利的简单语言，用以表达双线性型理论的结果，而后者的进展与不变量理论的进展齐头并进。\(^{24}\)凡涉及双线性型的秩的概念以及这些型与线性变换之间的关系的一切内容，都由弗罗贝尼乌斯的著作（{[}119{]}，第 I 卷，第 343–405 页）得到了最终的澄清。也是多亏弗罗贝尼乌斯，我们才有了一个交错型在自由 Z-模上的标准表达式（{[}119{]}，第 I 卷，第 482–544 页）；不过，斜对称行列式早在本世纪初已随普法夫出现，联系着把微分形式化为规范形的问题；雅可比在 1827 年重新研究了这个问题（{[}171{]}，第 IV 卷，第 17–29 页），他知道奇数阶的斜对称行列式为零，而且正是他构造出普法夫式的表达式并证明它是偶数阶斜对称行列式的一个因子；但他没有意识到后者是普法夫式的平方，这一点直到 1849 年才由凯莱建立（{[}58{]}，第 I 卷，第 410–413 页）。与一个二次型相伴的对称双线性型的概念是“配极化”过程的最简单情形，而配极化是不变量理论的基本工具之一。以“标量积”的名义，这一概念将享有巨大的昌盛：先是在“向量演算”的普及者那里，然后从 20 世纪起，得益于希尔伯特空间理论带给它的未曾预料到的推广（见第 212 页）。也正是后一种理论将使算子的伴随这一概念显露出来（在此之前，它几乎只在线性微分方程理论中露过面，而在张量演算中，它则表现为度量张量的指挥下协变指标与逆变指标的轮舞）；最终赋予埃尔米特型概念以其全部个性的也正是这一理论，该概念最初由埃尔米特于 1853 年在其算术研究中引入（{[}159{]}，第 I 卷，第 237 页），但直到大约 1925 年以及复希尔伯特空间应用于量子理论之前，它一直处于数学大潮的边缘。

正交群和相似群——自 19 世纪中叶起就被清楚地设想并照此处理，而且已经成为二次型理论的核心——以及其他“经典”群（线性群、辛群和酉群）的研究，则占据了越来越重要的地位。这些群在李群理论与微分几何中、以及在二次型算术理论（例如见{[}285{]}和{[}100{]}）中所起的重要作用，我们在此只能提上一句；\(^{25}\)正是由于这一情况，也由于对偶概念向最多样问题的推广，在现代数学理论中才几乎没有哪一门不以这种或那种形式牵涉到双线性型。无论如何应当指出，正是（三维）旋转群的研究引导哈密顿发现了四元数{[}145 a{]}；这一发现由 W. 克利福德加以推广，他于 1876 年引入了以他的名字命名的那些代数，并证明它们是若干四元数代数、或四元数代数与一个二次扩张的张量积（{[}65{]}，第 266–276 页）。这些代数四年后被利普希茨重新发现{[}205 b{]}，他用它们给出了 n 个变量的正交变换的参数表示（推广了凯莱借助四元数理论就 n = 3（{[}58{]}，第 I 卷，第 123–126 页）和 n = 4（第 II 卷，第 202–215 页）得到的那些表示），它们以及由之派生的“旋量”概念（{[}52 b{]}和{[}62 a{]}）也因其在量子理论中的应用而在现代享有巨大的流行。

最后，我们还须谈一谈思想的演变，它导致半双线性型理论中关于标量环的一切限制几乎全部被抛弃——这是整个现代代数中共同的倾向，但它在这里也许比在别处显现得更早。我们已经指出过复数域上的几何的富有成效的引入（此外，在整个 19 世纪，这种几何与实几何之间一直存在一种持久的、有时是危险的混淆）；这方面的清晰特别来自 19 世纪末关于几何基础的公理化研究{[}163 c{]}。在这类研究中，希尔伯特及其追随者在考察各条公理之间的关系时，被引去构造适当的反例，其中的“基域”（交换的或非交换的）具有或多或少的病态性质，他们由此使数学家们习惯于全新类型的“几何”。从分析的观点看，伽罗瓦已经考虑过系数和变量取值于有限素域的线性变换（{[}123{]}，第 145 页）；在发展这些思想时，若尔当{[}174 a{]}很自然地被引导去考虑这些域上的经典群，这些群的介入显现在各种各样的数学领域中。迪克森在大约 1900 年把若尔当的研究推广到一切有限域，而更近一些，人们已经看到，若尔当–迪克森理论的一大部分可以推广到绝对任意的“基域”上；这本质上应归功于迷向向量的一般性质和维特定理，它们在经典情形是平凡的，直到 1936 年才就任意基域建立起来{[}337 a{]}。\(^{26}\)

但是，在把半双线性型的研究推向不断增大的“抽象”的同时，人们发现，把二维或三维情形中出自经典几何的术语原样保留下来，并把它推广到 n 维情形乃至无穷维空间，是极富启发性的。经典几何就这样卸下了它作为一门自治的、有生命的科学的角色，被重塑为当代数学的一种普遍语言，其灵活性和有用性都是无与伦比的。

